% =====================================================================
%  Paper empírico — versión standalone para envío a revista
%  Fragilidad bancaria sistémica, riesgo de cola del crecimiento y
%  spread soberano en economías emergentes
%
%  Reutiliza el cuerpo del capítulo único de la tesis (../../tesis/investigacion.tex)
%  con remapeo de \chapter->\section y los entornos teoremáticos redefinidos.
%  Revistas objetivo (español): Latin American Journal of Central Banking,
%  Economía Chilena (BCCh), Estudios de Economía (U. de Chile), El Trimestre
%  Económico.
% =====================================================================
\documentclass[12pt,letterpaper]{article}
\usepackage[spanish,es-tabla]{babel}
\usepackage[utf8]{inputenc}
\usepackage[T1]{fontenc}
\usepackage{amsmath,amssymb,amsthm}
\usepackage{booktabs}
\usepackage{longtable}
\usepackage{array}
\usepackage{pdflscape}
\usepackage{graphicx}
\usepackage{tikz}
\usetikzlibrary{arrows.meta,positioning}
\newcommand{\tst}[1]{\ensuremath{\scriptstyle(#1)}}
\usepackage{setspace}
\usepackage[margin=1in]{geometry}
\usepackage[hidelinks]{hyperref}
\usepackage{natbib}
\bibliographystyle{apalike}

\graphicspath{{figuras/}}
\setlength{\emergencystretch}{3em}
\onehalfspacing

% --- macros de notación compartidas con la tesis ---
\newcommand{\JLoss}{\mathrm{JLoss}}
\newcommand{\Spread}{\mathrm{spread}}
\newcommand{\GaR}{\mathrm{GaR}}
\newcommand{\PD}{\mathrm{PD}}
\newcommand{\ES}{\mathrm{ES}}
\newcommand{\HHI}{\mathrm{HHI}}
\newcommand{\Var}{\mathrm{Var}}
\newcommand{\dd}{\,\mathrm{d}}
\newcommand{\E}{\mathbb{E}}
\newcommand{\1}{\mathbf{1}}
\providecommand{\R}{\mathbb{R}}
\let\epsilon\varepsilon

% --- entornos teoremáticos (sin contador [chapter]) ---
\theoremstyle{plain}
\newtheorem{proposition}{Proposición}
\newtheorem{lemma}{Lema}
\theoremstyle{definition}
\newtheorem{definition}{Definición}
\newtheorem{assumption}{Supuesto}
\newtheorem{remark}{Observación}

% --- remapeos para reusar el archivo de capítulo sin editarlo ---
\let\chapter\section
% El capítulo único (investigacion.tex) ya no referencia otros capítulos de la
% tesis (esa arista se eliminó); solo cita su propio anexo de datos, que este
% envío standalone no incluye -- se neutraliza a texto plano en su lugar.
\makeatletter
\newcommand{\definelabel}[2]{\expandafter\gdef\csname r@#1\endcsname{{#2}{}{}{}{}}}
\makeatother
\definelabel{chap:anexoA}{el anexo de procedencia de datos de la tesis}

\title{Fragilidad bancaria sistémica, riesgo de cola del crecimiento y
\textit{spread} soberano en economías emergentes}
\author{Mauricio Valenzuela Corvalán\\[2pt]
\small Magíster en Economía Aplicada, Universidad de Chile}
\date{\today}

\begin{document}
\maketitle

\begin{abstract}
\noindent
Este trabajo evalúa si la fragilidad sistémica del sector bancario y el riesgo de la
cola izquierda del crecimiento determinan de forma \emph{complementaria} —no aditiva—
el riesgo de crédito soberano en economías emergentes. La fragilidad se mide con
\textit{Joint Loss} ($\JLoss$), agregación de punto de silla de probabilidades de
incumplimiento tipo Merton sobre 113 bancos extraídos homogéneamente de Bloomberg; el
riesgo de cola del crecimiento, con \textit{Growth-at-Risk} ($\GaR$) por regresión
cuantílica de panel; el \textit{spread} soberano, por el EMBI Global Diversified de
J.P.\ Morgan. Sobre un panel único de trece economías emergentes (2004--2026, $N=721$),
los canales de nivel están bien identificados ($\hat\beta_1>0$ para $\JLoss$;
coeficiente positivo de $D\equiv-\GaR$). La interacción $\JLoss\times D$ tiene el signo
predicho por el mecanismo de \textit{doom loop} pero no es significativa sobre la
muestra completa ($\hat\beta_3=+0{,}16$, $p=0{,}26$). \textbf{Sí lo es} en dos
dimensiones: transversal, en el núcleo de once economías emergentes de financiamiento
externo ($\hat\beta_3=+0{,}47$, $p=0{,}023$); y temporal —interactuando la
complementariedad con un vector de crisis sin excluir observaciones, se \textbf{anula}
bajo respaldo oficial masivo (crisis financiera global, pandemia) y \textbf{sobrevive}
en el estrés emergente de 2015--2016 sin ese respaldo ($p<0{,}001$), un test de
falsación directo del mecanismo. El resultado es robusto a tratar $\JLoss$ y $\GaR$
como regresores generados y a la inferencia por \textit{wild cluster bootstrap}; la
identificación causal del canal de nivel, en cambio, no queda cerrada por variables
instrumentales. La complementariedad fragilidad--crecimiento sobre el riesgo soberano
es, así, un fenómeno \textbf{condicional} a la exposición al financiamiento externo y
al régimen de respaldo oficial, no una regularidad universal.\\[4pt]
\noindent\textbf{Palabras clave:} riesgo soberano, fragilidad bancaria, \textit{Growth-at-Risk},
\textit{doom loop}, crisis financiera, economías emergentes.\quad\textbf{JEL:} F34, G01,
G21, G15, H63, E44.
\end{abstract}

\input{../../tesis/investigacion.tex}

\end{document}
